% Root and DE independent audits passed on 2026-11-15.
\begin{statement}{C.34}{Proved}{Odd-primary arithmetic-trivial characters at all sufficiently large primes}
Let $V$ be the two-dimensional standard representation of $S_3$.
For $g\in S_3$, put
\[
 L_n(g)=\frac1n\sum_{d\mid n}\mu(d)
                   \operatorname{tr}(g^d\mid V)^{n/d}.
\]
Write $1,\tau,\rho$ for the identity, a transposition and a
three-cycle, and define, for $n\geq2$,
\[
\begin{split}
 a_n&=\frac{2L_n(1)-L_{n+1}(1)-2L_n(\rho)
                 -3L_{n+1}(\tau)-2L_{n+1}(\rho)}6,\\
 b_n&=\sum_{\substack{i+j=n\\2\leq i<j}}a_i a_j
              +\begin{cases}\binom{a_{n/2}}2,&2\mid n,\\0,&2\nmid n,\end{cases}\\
 d_n&=\max(0,a_n-b_n).
\end{split}                                                     \tag{C34.1}
\]
Let $Q$ be the full gentle group modulo the closed normal closure
of the absolute Galois image.  For every prime $p\geq5$,
\[
 \dim_{\mathbb F_p}\operatorname{Hom}_{\mathrm{cts}}(Q,\mathbb F_p)
                       \geq d_{p-1}.                            \tag{C34.2}
\]
Every displayed character can be chosen with its nonzero values
witnessed on global multiplier one.  In particular,
\[
\begin{array}{c|rrrrrr}
 p&13&17&19&23&29&31\\ \hline
 d_{p-1}&1&37&143&1805&77519&273536 .
\end{array}                                                     \tag{C34.3}
\]
The lower bound is positive and tends to infinity for every
sufficiently large prime $p$.  There is a continuous surjection
\[
 Q^{\mathrm{ab}}\longrightarrow
 \left(\prod_{j\geq1}C_2\right)
       \times\prod_{\substack{p\geq5\\p\ {\rm prime}}}
                         (\mathbb Z/p\mathbb Z)^{d_{p-1}} .
                                                               \tag{C34.4}
\]
This determines a large quotient, not the full abelianization.
No assertion is made here that the odd-primary characters
vanish on the pentagon group.
\end{statement}
\begin{proof}
We give an independent character computation of the dimensions
used in C.33.  Work first over $\mathbb Q$.
The free Lie algebra on $V$ has degree-$n$ piece $L_n(V)$.
A homogeneous tangential derivation is specified by three
conjugators $q_x,q_y,q_z\in L_n(V)$, with $x+y+z=0$.
For $n\geq2$ these conjugators have no ambiguity, because
the homogeneous peripheral centralizer in that degree is zero.
The special condition is
\[
            [x,q_x]+[y,q_y]+[z,q_z]=0 .
\]
Inner derivations are the diagonal triples.  The permutation
representation on the three conjugator positions modulo its
diagonal is $V$.  Bracketing with $x,y$ is onto $L_{n+1}(V)$
by Jacobi induction.  There is therefore an exact sequence
of $S_3$-representations
\[
 0\longrightarrow \mathfrak s_n
 \longrightarrow V\otimes L_n(V)
 \longrightarrow L_{n+1}(V)\longrightarrow0,                    \tag{C34.5}
\]
where $\mathfrak s_n$ is the space of special outer derivations.
Its $S_3$ invariants are exactly $S_n(\mathbb Q)$ of (C33.1).

For completeness, the trace formula in the statement follows
directly from the tensor algebra and its ordered Lie-monomial
basis.  The trace series of the tensor algebra at $g$ is
$(1-\operatorname{tr}(g\mid V)t)^{-1}$.
The Poincare--Birkhoff--Witt decomposition into symmetric
powers of the $L_n(V)$, followed by logarithms, gives
\[
 \operatorname{tr}(g\mid V)^m
       =\sum_{n\mid m}n\,\operatorname{tr}(g^{m/n}\mid L_n(V)).
\]
Mobius inversion gives $L_n(g)=\operatorname{tr}(g\mid L_n(V))$.
Since the three values of $\operatorname{tr}(g\mid V)$ are
$2,0,-1$, taking invariants in (C34.5) gives exactly $a_n$.
Thus $a_n$ is the dimension of the actual degree-$n$
gentle Lie space.
There are at most $b_n$ independent degree-$n$ brackets:
for $i<j$ there are $a_i a_j$ possible pairs, and for $i=j$
there are $\binom{a_i}{2}$.  Nested brackets are already
included, since their two outer factors are elements of
the same graded Lie algebra.  We have not assumed that
all these pairs are independent.

Fix $p\geq5$ and apply C.33's actual-image argument with
$13$ replaced by $p$ and top conjugator degree $p-1$.
The complete integral peripheral correction and coprime
$S_3$ argument give
\[
       \mathcal U_n/\mathcal U_{n+1}=S_n(\mathbb Z_p)
              \qquad(2\leq n\leq p-1).                          \tag{C34.6}
\]
These are free finite-rank $\mathbb Z_p$ modules, of ranks $a_n$.
The zero degree-two space again gives $\mathcal U=\mathcal U_3$.
On the Magnus algebra modulo augmentation degree $p+1$,
the operator $u-1$ raises degree by at least three.
The resulting actual quotient
$U=\mathcal U/\mathcal U_p$ has nilpotency class at most
$\lfloor(p-1)/3\rfloor<p$ and is a finitely generated
torsion-free pro-$p$ group.  Lazard logarithms, their
filtration and the leading Ihara bracket are integral.
This uses the small class of the automorphism group,
not a class-$p$ correspondence for the underlying free group.

A tame order-$(p-1)$ element in the local Galois image
has distinct eigencharacters on every possible grade
$3,\ldots,p-1$.  Its averaging idempotents split the
actual filtered lattice as in C.33.  The invariant
eigensummand is precisely the top piece of rank $a_{p-1}$.
Its intersection with the Lie bracket ideal has rank
at most $b_{p-1}$.  Hence the lattice of integral
functionals on that top piece which annihilate brackets
has rank at least $d_{p-1}$.
This lattice is saturated in the dual of the top piece:
if $p\ell$ annihilates every bracket, so does $\ell$,
because $\mathbb Z_p$ is torsion-free.
Consequently a basis supplies at least $d_{p-1}$
linearly independent reductions modulo $p$.
This saturation argument is what permits reduction
modulo $p$; rational dimension alone would not suffice.

Extending each functional by zero on the other tame
eigensummands gives an ordinary homomorphism
$U\to\mathbb F_p$.  It is invariant under the whole
local gentle image: conjugation on the abelianization
factors through the multiplier and acts on the top
summand by $\lambda_p^{p-1}=1$ modulo $p$.
The vanishing of
$H^2_{\mathrm{cts}}(\mathbb Z_p^\times,\mathbb F_p)$
extends each character to the full local image.
Its Galois restriction is unramified away $p$.
Kronecker--Weber makes that restriction a scalar multiple
of
\[
               c_p(\lambda)=
                    \frac{\lambda_p^{p-1}-1}{p}\pmod p .
                                                               \tag{C34.7}
\]
Subtracting this scalar multiple makes the character
arithmetic-trivial.  The subtraction vanishes on
global multiplier one, where (C34.6) and P.4 provide
actual complete witnesses for the independent values.
Thus their independence is retained.  This proves
(C34.2), with no assumption about generation of the
arithmetic Lie algebra or its integral image.

We prove the assertion about sufficiently large primes.
Put $l_n=L_n(1)$.  The trace formula gives
\[
 l_n=\frac{2^n}{n}+O(2^{n/2}),\qquad
 a_n=\frac{2^n}{3n(n+1)}+O(2^{n/2})
          \sim\frac{2^n}{3n^2}.                                \tag{C34.8}
\]
For the nonidentity classes the largest exponential
terms have sizes $2^{n/2}$ and $2^{n/3}$, so they are
included in the displayed error.
In particular $r_n=a_n2^{-n}\sim(3n^2)^{-1}$.
The sum $S=\sum_{n\geq2}r_n$ is finite and small.
Indeed (C34.5) gives $a_n\leq2l_n-l_{n+1}$, and
\[
 \sum_{n=2}^{N}\frac{2l_n-l_{n+1}}{2^n}
              =\frac12-\frac{l_{N+1}}{2^N}
                 \longrightarrow\frac12.
                                                               \tag{C34.9}
\]
At $n=5$ this inequality is strict by two dimensions,
so in fact $S\leq7/16<1/2$.

The elementary convolution estimate for
$r_n\sim(3n^2)^{-1}$ is
\[
 \frac{\sum_i r_i r_{n-i}}{r_n}\longrightarrow2S.                \tag{C34.10}
\]
To verify it, first retain only $i\leq M$ and $n-i\leq M$.
Their ratio tends to $2\sum_{i\leq M}r_i$.
On the remaining half interval $M<i\leq n/2$,
$r_{n-i}/r_n$ is uniformly bounded, while
$\sum_{i>M}r_i=O(1/M)$.  The other half is identical;
then let $M$ tend to infinity.
The diagonal correction in $b_n$ is exponentially
smaller than $a_n$.  Therefore
\[
 \frac{b_n}{a_n}\longrightarrow S<\frac12,\qquad
 d_n\sim(1-S)\frac{2^n}{3n^2}.                                  \tag{C34.11}
\]
It follows that $d_{p-1}$ is positive and unbounded
for every sufficiently large prime.
The exact formula program
\texttt{pbc\_gentle\_dimension\_formula.py} verifies
(C34.3), agrees with all independently computed
dimensions through degree twelve, and records additional
counts through degree $300$.  Those finite computations
illustrate (C34.11); they are not the proof of its
eventual assertion.

For each prime choose the independent characters just
constructed.  They jointly surject onto the indicated
finite elementary abelian $p$-group.  Any finite product
over distinct primes is still a quotient: a subdirect
product of finite abelian groups of pairwise coprime
orders is their product.  The same argument includes
arbitrary finite subsets of C.30's independent
quadratic characters.  Compactness now gives
surjectivity onto the full product in (C34.4).
\end{proof}
